%%%%%%%%%%%%%%%%%%%%%%%%%%%%%%%%%%%%%%%%%
% Classicthesis Typographic Thesis
% LaTeX Template
% Version 1.4 (1/1/16)
%
% This template has been downloaded from:
% http://www.LaTeXTemplates.com
%
% Original author:
% André Miede (http://www.miede.de) with commenting modifications by:
% Vel (vel@LaTeXTemplates.com)
%
% License:
% GNU General Public License (v2)
%
% General Tips:
% 1) Make sure to edit the classicthesis-config.file
% 2) New enumeration (A., B., C., etc in small caps): \begin{aenumerate} \end{aenumerate}
% 3) For margin notes: \marginpar or \graffito{}
% 4) Do not use bold fonts in this style, it is designed around them
% 5) Use tables as in the examples
% 6) See classicthesis-preamble.sty for useful commands
%
%%%%%%%%%%%%%%%%%%%%%%%%%%%%%%%%%%%%%%%%%

%----------------------------------------------------------------------------------------
%	PACKAGES AND OTHER DOCUMENT CONFIGURATIONS
%----------------------------------------------------------------------------------------

% Os parâmetros "oneside" e "openany" fazem com que não hajam páginas em branco na impressão 
\documentclass[
		oneside,openany,titlepage,numbers=noenddot,headinclude,%headlines,
	 	footinclude=true,
		dottedtoc, % Make page numbers in the table of contents flushed right with dots leading to them
		BCOR=5mm,paper=a4,fontsize=11pt, % Binding correction, paper type and font size
		portuguese, % Languages, change this to your language(s)
		]{scrreprt} 
                
% Includes the file which contains all the document configurations and packages - make sure to edit this file
\input{classicthesis-config}

% Espaçamento entre linhas
\usepackage{setspace}

%\usepackage[document]{ragged2e}

% Caracteres especiais ao longo do texto
%\usepackage[utf]{inputenc}

\usepackage{geometry}
 \geometry{
 a4paper,
 %total={170mm,257mm},
 left=30mm,
 right=20mm, % O correto seria 20mm
 top=30mm,
 bottom=30mm
}

% Mathematical tools
\usepackage{siunitx}
\usepackage{amssymb}
\usepackage{mathtools}
\usepackage{amsmath}
\usepackage{relsize}
\usepackage{gensymb}

% Timeline draw
\usepackage{charter}
\usepackage{environ}
\usepackage{tikz}
\usetikzlibrary{calc,matrix}

%Sistemas de equações
\usepackage{systeme}

% Rotacionar tabelas
\usepackage{adjustbox}
\usepackage{rotating}

% Tabelas
\usepackage{booktabs}

%Legendas
\usepackage[labelfont=bf,textfont=it]{caption}

% Permite o uso de estruturas enumerativas
\usepackage{enumitem}

% Float figures
\usepackage{float}

% Caracteres unicode
\usepackage[utf8]{inputenc}

%\usepackage[sectionbib]{chapterbib}
\addbibresource{Bibliography.bib} % The file housing your bibliography
%\usepackage[num]{abntex2cite}
%\renewcommand{\bibname}{Bibliography}
%\addbibresource[label=ownpubs]{Self_Publications.bib} % Uncomment for optional self-publications
%\citeoption{abnt-options4}

%\hyphenation{Put special hyphenation here}

% Comando para alterar a fonte de todo o trabalho
%\renewcommand{\familydefault}{\sfdefault}
\renewcommand{\familydefault}{cmr}

% Identação na primeira linha
\setlength{\parindent}{3em}

\usepackage{titlesec}
\titleformat{\chapter}[hang]{\huge}{\thechapter}{1em}{}
\titlespacing{\chapter}{0pt}{0pt}{1cm}

\begin{document}

\frenchspacing % Reduces space after periods to make text more compact

\raggedbottom % Makes all pages the height of the text on that page

\selectlanguage{portuguese} % Select your default language - e.g. american or ngerman

%\renewcommand*{\bibname}{new name} % Uncomment to change the name of the bibliography
%\setbibpreamble{} % Uncomment to include a preamble to the bibliography - some text before the reference list starts

\pagenumbering{roman} % Roman page numbering prior to the start of the thesis content (i, ii, iii, etc)

\pagestyle{plain} % Suppress headers for the pre-content pages

%----------------------------------------------------------------------------------------
%	PRE-CONTENT THESIS PAGES
%----------------------------------------------------------------------------------------

\include{FrontBackMatter/Titlepage} % Main title page 

\include{FrontBackMatter/Backcover} % Main title page

\include{FrontBackMatter/Dedication} % Dedication page

\include{FrontBackMatter/Abstract} % Abstract page

\include{FrontBackMatter/Acknowledgments} % Acknowledgements page

\pagestyle{scrheadings} % Show chapter titles as headings

\include{FrontBackMatter/Contents} % Contents, list of figures/tables/listings and acronyms

\pagenumbering{arabic} % Arabic page numbering for thesis content (1, 2, 3, etc)
%\setcounter{page}{90} % Uncomment to manually start the page counter at an arbitrary value (for example if you wish to count the pre-content pages in the page count)

%----------------------------------------------------------------------------------------
%	THESIS CONTENT - CHAPTERS
%----------------------------------------------------------------------------------------
%\begin{refsection}

\include{Chapters/Introduction} % Introdução
\include{Chapters/Objectives} % Objetivos
\include{Chapters/Theory} % Revisão bibliográfica
\include{Chapters/Methodology} % Metodologia
\include{Chapters/Results} % Revisão bibliográfica
\include{Chapters/Conclusions} % Conclusões e perspectivas
%\printbibliography
%\include{FrontBackMatter/Bibliography}

%\end{refsection}
%------------------------------------------------

\appendix
\ctparttext{Ao longo dos Apêndices são apresentados os resultados}
\part{Apêndice} % New part of the thesis for the appendix

\begin{refsection}
\include{Chapters/Appendix} % Appendix A
\end{refsection}

%----------------------------------------------------------------------------------------
%	POST-CONTENT THESIS PAGES
%----------------------------------------------------------------------------------------
\include{FrontBackMatter/Bibliography} % Bibliography
\include{FrontBackMatter/Timeline} % Timeline

%----------------------------------------------------------------------------------------

\end{document}
